\section{Pinning zeros and transporting weighted observables}
\label{sec:pinning}

We compare zero modes at fluxes differing by one or three.  Multiplication
by a zero at the south pole gives the algebraic comparison, but is not
unitary.  We construct a unitary comparison whose effect on a local
observable is small away from that pole.  For a convex orbital weight we
also retain a sign in this comparison.  That sign will be used in the first
charge estimate in Section~\ref{sec:charge}.

Fix $d\in\{1,3\}$ and put $t=q-d$.  The identification of orbital labels
$0,\ldots,t$ embeds $U_t$ isometrically into $U_q$; write $\iota_{q,d}$ for
its exterior-algebra extension.  Its image has the last $d$ orbitals empty.
Set
\begin{equation}
 g_{q,d}(j)=\frac12\log\frac{\binom{t}{j}}{\binom qj}
 =\frac12\sum_{r=0}^{d-1}
       \log\frac{q-j-r}{q-r},\qquad 0\le j\le t.
 \label{eq:pinning-log-weight}
\end{equation}
Multiplication of a polynomial by the $d$ south-pole zero factors, one
factor for each particle, is $\iota_{q,d}\exp D_t(g_{q,d})$, up to a scalar
on each number sector.  In particular, the subspace of $Z_{n,q}$ having
these zeros is
\begin{equation}
 \iota_{q,d}\exp D_t(g_{q,d}) Z_{n,t}.
 \label{eq:pinning-kernel-identification}
\end{equation}
Throughout this section, $m$, $u=1+m$, $v=1+q-m$, and
$w=(uv/(q+1))^{1/2}$ use the larger flux $q$.  Replacing $q$ by $t$ in
these size functions changes them by bounded factors.  For every integer
$\ell\ge1$, the smooth extension of Equation~\eqref{eq:pinning-log-weight}
satisfies, for $0\le m\le t$,
\begin{equation}
 g_{q,d}^{(\ell)}(m)
 =-\frac{(\ell-1)!}{2}\sum_{r=0}^{d-1}(q-m-r)^{-\ell},
 \qquad |g_{q,d}^{(\ell)}(m)|\le C_\ell v^{-\ell},
 \qquad -g_{q,d}''(m)\asymp v^{-2}.
 \label{eq:pinning-derivatives}
\end{equation}
Here and below constants can depend on $d$ and on a stated differentiation
or localization order, but not on $q$, $n$, or $m$.

\subsection{A uniformly gapped path}

For $1\le b\le2t-1$, let $B_{t,b}$ be the normalized pair annihilator
whose coefficient at $i<j$, $i+j=b$, is proportional to
$(j-i)\sqrt{\binom ti\binom tj}$.  For $0\le s\le1$ define
\begin{equation}
 \begin{aligned}
 B_{q,d,b}(s)&=\sum_{\substack{0\le i<j\le t\\i+j=b}}
             a_{q,d,b,s}(i,j)c_jc_i,\\
 a_{q,d,b,s}(i,j)&=
 \frac{a_{t,b}(i,j)e^{-s(g_{q,d}(i)+g_{q,d}(j))}}
 {\left(\sum_{k<l,\ k+l=b}|a_{t,b}(k,l)|^2
          e^{-2s(g_{q,d}(k)+g_{q,d}(l))}\right)^{1/2}}.
 \end{aligned}
 \label{eq:pinning-rows}
\end{equation}
and
\[
 H_{q,d}(s)=\sum_b B_{q,d,b}(s)^*B_{q,d,b}(s),\qquad
 P_{q,d}(s)=\text{the orthogonal projection onto }\ker H_{q,d}(s).
\]
All these operators act on $\mathcal F_t$.  Direct conjugation of each
annihilator gives
\begin{equation}
 \ker H_{q,d}(s)=e^{sD_t(g_{q,d})}\ker H_t.
 \label{eq:pinning-path-kernel}
\end{equation}
The normalization in Equation~\eqref{eq:pinning-rows} only multiplies each
row by a positive scalar and therefore does not affect this identity.

\begin{lemma}[Gap along the pinning path]
\label{pinning:uniform-gap}
There exist $q_{\mathrm{pin}}$ and $\gamma_{\mathrm{pin}}>0$ such that
\[
 H_{q,d}(s)\ge\gamma_{\mathrm{pin}}(1-P_{q,d}(s))
 \quad(q\ge q_{\mathrm{pin}},\ d\in\{1,3\},\ 0\le s\le1)
\]
on the entire Fock space $\mathcal F_{q-d}$.
\end{lemma}

\begin{proof}
Fix an allowed unperturbed Fock gap $\gamma>0$.  We compare columns of
annihilators, rather than Hamiltonians.  This avoids a factor equal to the
number of particles.

Choose a smooth function $\ell_B(x)$ on $x\ge1$, constant for $x\le B$,
equal to $\log x$ for $x\ge2B$, and with
$|\ell_B^{(k)}(x)|\le C_k(B+x)^{-k}$ for $k\ge1$.  Replace each logarithm
$\log(q-j-r)$ in Equation~\eqref{eq:pinning-log-weight} by
$\ell_B(q-j-r)$, retaining its constant denominator term, and denote the
result by $\bar g$.  Let $\bar B_b(s)$ be the corresponding normalized
rows.  At $m=b/2$, the pair-row moment bounds of
Lemma~\ref{locality:row-symbols} imply
\begin{equation}
 \|\bar B_b(s)-B_{t,b}\|
 \le C\sqrt w\,\frac{w^2}{(v+B)^2}.
 \label{eq:pinning-rounded-row-error}
\end{equation}
Here is the normalization detail in this application.  Subtract the affine
Taylor polynomial of $\bar g$ at $m$ from each of its two arguments.  The
subtracted contribution is constant on $i+j=b$ and cancels in row
normalization.  Write $z=|i-m|/w$.  Taylor's formula, the derivative bounds
on $\ell_B$, and $w^2\le v$ give a bound
$Cw^2(v+B)^{-2}(1+z)^K$ for the remaining exponent; exponentiating costs
at most another fixed power of $1+z$.  To check this also outside
$|i-m|\le(v+B)/2$, observe that this exterior region has
$z\ge c\sqrt v$ whenever $v\ge B$, whereas logarithmic differences are
bounded by $C\log(1+v/B)$.  If $v<B$, all the logarithms have bounded
scaled derivatives on the whole row.  Thus the elementary inequality
$|e^x-1|\le |x|e^{|x|}$ and the row moment estimates give both an
$\ell^2$ coefficient error $Cw^2/(v+B)^2$ and an $\ell^1$ error
$C\sqrt w\,w^2/(v+B)^2$.  The normalization factor differs from one by
the first of these errors.  Since $w^2/(v+B)^2\le C/B$, it is bounded
away from zero when $B$ is large.  Finally
$\|\sum e_{ij}c_jc_i\|\le\sum|e_{ij}|$ proves
Equation~\eqref{eq:pinning-rounded-row-error}.

Let $\mathcal C_0\psi=(B_{t,b}\psi)_b$ and
$\bar{\mathcal C}_s\psi=(\bar B_b(s)\psi)_b$.  These are operators from
$\mathcal F_t$ to a direct sum of copies of $\mathcal F_t$.  Equation~
\eqref{eq:pinning-rounded-row-error} gives
\begin{equation}
 \|\bar{\mathcal C}_s-\mathcal C_0\|^2
 \le C\sum_b\frac{w_b^5}{(v_b+B)^4}
 \le C\int_0^\infty\frac{(1+x)^{5/2}}{(B+1+x)^4}\,dx
 \le C B^{-1/2}.
 \label{eq:pinning-column-error}
\end{equation}
The middle inequality uses $w_b^2\le v_b$ and the spacing $1/2$ of
$v_b$.  Consequently the operator-column error is at most $CB^{-1/4}$.
Choose $B$, independent of $q$ and $s$, so that this is at most
$\sqrt\gamma/2$.  Equation~\eqref{eq:pinning-path-kernel}, also valid
for $\bar g$, shows that the two columns have the same nullity in each
number sector.  The min--max principle for singular values and the
supplied Fock gap therefore give
\[
 \|\bar{\mathcal C}_s\psi\|
 \ge\frac{\sqrt\gamma}{2}
       \operatorname{dist}(\psi,\ker\bar{\mathcal C}_s).
\]

It remains to remove the rounding.  Put $\delta=g_{q,d}-\bar g$ and
$S_s=e^{sD_t(\delta)}$.  At most $2B+d$ orbital entries of $\delta$ are
nonzero, and each is bounded in modulus by $C_d\log(2B)$.  Thus
\[
 \|S_s\|\|S_s^{-1}\|
 \le\exp\Bigl(\sum_j|\delta(j)|\Bigr)=:K_B<\infty
\]
uniformly on the entire Fock space.  If $\mathcal C_s$ is the column of
the unrounded normalized rows, row normalization gives
\[
 \mathcal C_s=R_s(\mathbb 1\otimes S_s)
                 \bar{\mathcal C}_s S_s^{-1},
 \qquad \|R_s\|+\|R_s^{-1}\|\le C_B,
\]
where $R_s$ is diagonal in the row index.  Since
$\ker\mathcal C_s=S_s\ker\bar{\mathcal C}_s$, the preceding distance
estimate implies
\[
 \|\mathcal C_s\psi\|
 \ge \frac{\sqrt\gamma}{2C_BK_B}
             \operatorname{dist}(\psi,\ker\mathcal C_s).
\]
Squaring proves the assertion.  The flux threshold only has to ensure
$t\ge q_\gamma$ and $q\ge4B+d$.
\end{proof}

\subsection{Local transport and comparison of consecutive fluxes}

For the remainder of this section all compared path fluxes are assumed
to be at least $q_{\mathrm{pin}}$.  Taking $q\ge q_{\mathrm{pin}}+6$
ensures this for the comparisons of size at most three and the growth
steps used below.  Larger fixed comparison ranges require only a
corresponding fixed increase in this margin.

We compare different orbital dimensions by a fixed convention.  If
$t\le t'$, identify orbitals with the same labels and extend an even
operator $A$ on $\mathcal F_t$ to $A\otimes I$ on $\mathcal F_{t'}$,
using the graded tensor decomposition over the additional modes.  Its original
matrix elements are recovered by putting those modes in the vacuum.
We then use the common slot frames of
Lemma~\ref{locality:common-frames}.  All operator differences below use
this padding convention.

The local rate of the transport will be
\begin{equation}
 \alpha_q(m)=\frac{w^2}{v^2}.
 \label{eq:pinning-alpha}
\end{equation}
In the north this rate is of order $(1+m)/q^2$; changing the flux by a
bounded amount will gain one more inverse power of $v$.

\begin{proposition}[Transport of the pinned spaces]
\label{pinning:generator}
There is a number-preserving and axially invariant unitary path
$U_{q,d}(s)$ on $\mathcal F_t$, starting at the identity, such that
\[
 U_{q,d}(s)\ker H_t=\ker H_{q,d}(s).
\]
Writing $\partial_sU_{q,d}(s)=Y_{q,d}(s)U_{q,d}(s)$, the anti-Hermitian
generator $Y_{q,d}(s)$ is a rapid interaction of size $C\alpha_q(m)$.
In common frames, a bounded change in $q$, with $d$ fixed, changes the
generator interaction by size $C\alpha_q(m)/v$.
\end{proposition}

\begin{proof}
We first prove uniform row regularity and its flux difference, then
estimate the affine-subtracted commutator that enters the spectral
filter.

\emph{Row regularity.}
The unnormalized coefficients in Equation~\eqref{eq:pinning-rows} are
\[
 (j-i)\Bigl[
  \binom ti^{1-s}\binom qi^s
  \binom tj^{1-s}\binom qj^s\Bigr]^{1/2}.
\]
They obey the pair-row bounds in Lemma~\ref{locality:row-symbols}
uniformly for $0\le s\le1$.  We give the normalization argument because
the normalization is now a sum with moving endpoints.  Write $z=i-m$.
Relative to the undeformed row, after canceling its midpoint value, the
additional factor is exactly
\begin{equation}
 E_s(m,z)=\prod_{r=0}^{d-1}
       \left(1-\frac{z^2}{(q-m-r)^2}\right)^{-s/2}.
 \label{eq:pinning-profile-factor}
\end{equation}
When $v$ is bounded, all allowed $z$ and $w$ are bounded and the estimates
follow from finite differences.  When $v$ is large and $|z|\le c v$,
the logarithm of this factor and its derivatives in scale-$w$ variables
are bounded by $C(w^2/v^2)$ times fixed polynomials in $|z|/w$.
The positive exponential factor can be absorbed into half the Gaussian
decay of the undeformed row, since its exponent is $O(z^2/v^2)$ while
that decay has exponent $-c z^2/w^2$ and $w^2\le v$.
For $|z|>c v$, the factor in
Equation~\eqref{eq:pinning-profile-factor} is at most $C v^{d/2}$,
whereas $|z|/w\ge c\sqrt v$.  The Gaussian tail therefore pays this
factor and every required inverse power of $w$.

To control normalization differences, first retain a smooth central
cutoff supported in $|z|\le c\min(u,v)$ and equal to one on a smaller
such interval.  Use the smooth log-gamma expression for the undeformed
row before normalization, and sum its square times $E_s^2$ and this
cutoff over all fixed integers $i$.  This is a smooth function of real
$m$, bounded below by $c w^3$ after extracting the common midpoint
binomial factor: the $O(w)$ indices with $|i-m|\asymp w$ contribute
$\asymp w^2$ each.  Its derivative of order $k$ is at most
$C_k w^{3-k}$, by the Gaussian moment bounds.  The discarded tail is
$O_D(w^{3-D})$ for every $D$, also after any fixed finite difference:
bound that difference by its finitely many values at adjacent centers,
where the scales are comparable.  The chain rule for the retained
positive normalization, followed by this tail estimate, proves the
full finite-difference estimates for the exactly normalized rows.
For bounded $w$ the normalization is bounded above and below directly.
Thus no derivative of a moving-endpoint sum is being assumed.
Lemma~\ref{locality:row-symbols} now proves that $H_{q,d}(s)$ is a rapid
interaction of bounded strength.

\emph{Flux differences.}
Use common windows for fluxes a bounded distance apart, as in
Lemma~\ref{locality:common-frames}.  With orbital labels fixed from the
north, their Hamiltonian interactions differ by size $C/v$.  Indeed
\[
 \log\binom{Q+1}{j}-\log\binom Qj
     =\log(Q+1)-\log(Q+1-j).
\]
For the path profiles, after canceling their common midpoint values,
the corresponding ratio is exactly
\[
 \left(1-\frac{z^2}{(t+1-m)^2}\right)^{-(1-s)/2}
 \left(1-\frac{z^2}{(q+1-m)^2}\right)^{-s/2}.
\]
Its difference from one has size $Cw^2/v^2\le C/v$ in the row-symbol
bounds.  The preceding central-cutoff normalization
argument, applied to the difference of the two profiles, retains this
extra factor in every finite-difference bound.  In using that argument,
the factor $w^2/v^2$ is retained throughout $|z|\le cv$, including the
part discarded by the smaller central cutoff.  Only on $|z|>cv$ is it
recovered from tail powers, using
$v^2/w^2\le C(|z|/w)^2$.  This distinction is necessary when $v$ is
much larger than $u$.  If an individual pair
coefficient is present only at the larger flux, its index excursion
satisfies $|i-m|\ge v-O(1)$.  For large $v$ its Gaussian tail pays both
the inverse of $w^2/v^2$ and every required symbol seminorm; for bounded
$v$ it is a bounded-scale estimate.  The entirely new rows lie in a
bounded south zone, where $1/v$ is bounded below.  This proves the
claimed common-frame interaction difference, including the changed
endpoint supports.  Repeating the one-flux identity handles any fixed
bounded change in flux.

\emph{Affine subtraction.}
For a smooth real
function $a$ on the orbital interval, set
$X_a=[H_{q,d}(s),D_t(a)]$.  Decompose the Hamiltonian into annuli by a
sum-of-squares partition $\sum_k\vartheta_k(b)^2=1$, of row-index width
comparable to $w_k$, and then use its Fourier wavepackets.  If $m_k$ is
the annulus center, put
\[
 a_k^\circ(j)=a(j)-a(m_k)-a'(m_k)(j-m_k).
\]
The affine contribution commutes with the annular sum, because every row
has a fixed pair index sum.  More explicitly, if $C_x$ is its Fourier
wavepacket and $C_x(a)$ is the same wavepacket with row coefficients
multiplied by $a_k^\circ(i)+a_k^\circ(j)$, then
\begin{equation}
 X_a=\int X_{a,x}\,d\mu(x),\qquad
 X_{a,x}=C_x^*C_x(a)-C_x(a)^*C_x.
 \label{eq:pinning-source-decomposition}
\end{equation}
The center measure has bounded mass in unit balls.  The wavepacket
normalization is $w_k^{-1/2}$ and its circle measure is
$w_k\,d\phi/(2\pi)$, so Fourier orthogonality proves
Equation~\eqref{eq:pinning-source-decomposition} exactly.
The weighted row bounds give size $Cw^2|a''(m)|$ whenever
Equation~\eqref{eq:affine-symbol-hypothesis} holds with
$\beta=w^2|a''(m)|$.  For $a=g_{q,d}$ this size is $C\alpha_q(m)$:
for every $k\ge2$,
\[
 w(j)^k|g_{q,d}^{(k)}(j)|
 \le C_k\frac{w(j)^2}{v(j)^2}
                \left(\frac{w(j)}{v(j)}\right)^{k-2}
 \le C_k\alpha_q(j),
\]
and $\alpha_q$ has tempered ratio exponent four, independently of $k$.
Upon changing $q$ by a bounded amount, the source interaction changes by
size $C\alpha_q/v$: differences of $g''$ gain $v^{-1}$, and the row
wavepackets change by $C/v$.  These assertions include every fixed rapid
localization order, since the estimates just used also apply after the
finite differences used in Fourier localization.

\emph{Spectral filter.}
We use the spectral-filter construction of quasi-adiabatic
continuation~\cite{HW,BMNS}, with locality supplied by Section~\ref{sec:locality}.
Choose a smooth odd real function $\eta$ that is zero on
$[-\gamma_{\mathrm{pin}}/2,\gamma_{\mathrm{pin}}/2]$ and equals
$\operatorname{sgn}(\omega)$ for $|\omega|\ge\gamma_{\mathrm{pin}}$.
Let $r(\omega)=\eta(\omega)/\omega$, with value zero near zero, and use
the spectral calculus of $\operatorname{ad}_{H}(A)=[H,A]$ to set
\[
 Y_{q,d}(s)=r(\operatorname{ad}_{H_{q,d}(s)})X_g.
\]
The function $r$ and each of its derivatives belong to $L^2(\mathbb R)$.
Its inverse Fourier transform therefore has every absolute polynomial
moment: use Cauchy--Schwarz on $|t|\le1$, and the $L^2$ norm of a
sufficiently high derivative of $r$ on $|t|>1$.  Lemma~\ref{locality:filters}
and Equation~\eqref{eq:pinning-source-decomposition} give rapid locality
with size $\alpha_q$.

The spectral multiplier of $D_t(g)$ is $\eta(\omega)$, so $Y^*=-Y$.
If $P=P_{q,d}(s)$ and $Q=1-P$, the gap gives
\[
 YP=QD_t(g)P,\qquad PY=-PD_t(g)Q,\qquad PYP=0.
\]
Differentiating Equation~\eqref{eq:pinning-path-kernel} gives
$\partial_sP=QD_t(g)P+PD_t(g)Q=[Y,P]$.  The unitary solution therefore
transports $P$.  Both $H_{q,d}(s)$ and $D_t(g)$ commute with number and
with $\sum_jj c_j^*c_j$, proving the two symmetries.

Finally compare two fluxes in common frames.  The change of $X_g$ has
size $\alpha_q/v$.  The Hamiltonian change has size $1/v$.  Inserting
the latter into the time evolution in the filter by Duhamel's formula
costs $\alpha_q/v$, by Lemma~\ref{locality:difference}.  All filters are
the same for every flux and every path parameter.  This proves the
difference assertion.
\end{proof}

The isometry used subsequently is
$\mathcal U_{q,d}=\iota_{q,d}U_{q,d}(1)$ restricted to $Z_{n,q-d}$.
By Equation~\eqref{eq:pinning-kernel-identification}, its image is exactly
the subspace with $d$ south-pole zeros.

\begin{corollary}[A fixed observable in the northern region]
\label{pinning:northern-observable}
Let $O$ be a neutral operator of norm at most one supported on $B_r(x)$,
where $r\ge2$ and $x$ has latitude index $m$.  There are fixed constants
$b,C_*$, and $c=8$, independent of every tail order below, such that,
if $q>Cr^b(1+m)$, the operator
\[
 \Delta_{q,d}(O)=U_{q,d}(1)^*OU_{q,d}(1)-O
\]
satisfies
\begin{align}
 \|\Delta_{q,d}(O)\|&\le C r^c\frac{1+m}{q^2},\label{eq:pinning-northern-norm}\\
 \inf_{A\in\mathfrak A(B_{r+R}(x))}\|\Delta_{q,d}(O)-A\|
 &\le C_D r^c\frac{1+m}{q^2}(1+R)^{-D}
 \quad(R\ge C_*(1+r)).\label{eq:pinning-northern-tail}
\end{align}
In particular, it admits a decomposition
\begin{equation}
 \Delta_{q,d}(O)=\sum_{\nu\ge0}D_\nu,\qquad
 D_\nu\in\mathfrak A(B_{C_*2^\nu r}(x)),\qquad
 \|D_\nu\|\le C_D r^c\frac{1+m}{q^2}2^{-\nu D}.
 \label{eq:pinning-northern-shells}
\end{equation}
For $|q-q'|\le3$, with the same $d$ and identical northern slots,
$\Delta_{q,d}(O)-\Delta_{q',d}(O)$ obeys all these estimates with
$q^{-2}$ replaced by $q^{-3}$.  In particular, for every fixed $B$,
the shell decomposition can be chosen so that
\begin{equation}
 \sum_{\nu\ge0}(C_*2^\nu r)^B\|D_\nu\|
 \le C_B r^{B+c}\frac{1+m}{q^2},
 \label{eq:pinning-northern-moment}
\end{equation}
and the same moment estimate holds for the flux difference with $q^{-3}$.
\end{corollary}

\begin{proof}
At the center the cutoff gives $v\asymp q$, so
$\alpha_q(m)\le C(1+m)/q^2$ and
$\alpha_q(m)/v(m)\le C(1+m)/q^3$.
Use the supported-observable bounds in
Lemma~\ref{locality:difference}, first for the self-adjoint path
generators $iY_{q,d}(s)$ and $0$.  Since
$\alpha_q=u/((q+1)v)$, its tempered ratio exponent can be taken to be
four: each of $u$ and $v^{-1}$ costs at most two powers of distance.
Equations~\eqref{eq:difference-core}--\eqref{eq:difference-tail} give
a core factor $r^6$ and every tail order beyond radius $Cr$.
Apply the same bounds directly to $iY_{q,d}(s)$ and
$iY_{q',d}(s)$ for the comparison.  Their difference has size
$\alpha_q/v=u/((q+1)v^2)$, whose ratio exponent is at most six, and
therefore costs $r^8$.  Thus $c=8$ suffices for both assertions.
Choose $b$ large enough that the initial ball uses identical northern
slots.  The common-frame generator comparison already includes the
padded south modes, so this same bound covers the endpoint changes.
Take the ball approximations by the normalized partial trace onto the
corresponding slot algebra.  This contraction has error at most twice
the best approximation error and preserves number grade.  Differences
at radii $C_*2^\nu r$ give one decomposition satisfying
Equation~\eqref{eq:pinning-northern-shells} for all tail orders,
after increasing $C_*$ by one fixed factor.  Choose the tail order
greater than $B+1$ and sum
the resulting geometric series to obtain
Equation~\eqref{eq:pinning-northern-moment}.
\end{proof}

\subsection{Convex weighted observables}

For fixed $0<p<2$ and $L\ge2$, put
\[
 f_{L,p}(j)=\left(\frac L{L+j}\right)^p,\qquad
 \beta_q(m)=w^2 f_{L,p}''(m),\qquad
 A_q(m)=\alpha_q(m)\beta_q(m).
\]
Constants in the next statements can depend on $p$, but are uniform in
$L\ge2$.  Direct calculation gives
\begin{equation}
 A_q(m)=\frac{u^2}{(q+1)^2}f_{L,p}''(m)
 \asymp\frac{f_{L,p}(m)}{(q+1)^2}
           \left(\frac{u}{L+u}\right)^2.
 \label{eq:pinning-weight-size}
\end{equation}
We use $d\mu$ for the center measure above; its integral of a radial
tempered size is bounded by a constant times $\int_0^t(\cdot)\,dm$.
These weights satisfy the precise symbol hypothesis needed for affine
subtraction.  For every $k\ge2$,
\[
 w(j)^k|f_{L,p}^{(k)}(j)|
 \le C_{k,p}\beta_q(j)
             \left(\frac{w(j)}{L+j}\right)^{k-2}
 \le C_{k,p}\beta_q(j).
\]
The ratios of $\beta_q(j)$ have a fixed polynomial bound in meridional
distance, independent of $k$ and $L\ge2$.  Thus
Equation~\eqref{eq:affine-symbol-hypothesis} applies with
$\beta=\beta_q(m)$ at each chosen center.

\begin{proposition}[Change of a weighted number observable]
\label{pinning:weighted-transport}
The operator
\[
 K_{q,d}(f)=U_{q,d}(1)^*D_t(f)U_{q,d}(1)-D_t(f)
\]
has a rapid interaction representation of size $CA_q(m)$.  Its
interaction difference under a bounded change in $q$, with $d$ fixed,
has size $CA_q(m)/v$ in common frames.  The pieces may be assigned to
their original annular centers throughout the construction.
\end{proposition}

\begin{proof}
Retain the annular decomposition of $Y$ used in its construction.
Each annular total commutes with number and axial angular momentum:
this is true of its source in Equation~\eqref{eq:pinning-source-decomposition}
and remains true under the time filter.  We may consequently replace
$f(j)$, in its commutator with this annular total, by
$f(j)-f(m_k)-f'(m_k)(j-m_k)$.  Lemma~\ref{locality:one-body-symbol}
bounds the resulting one-body interaction near the annulus by $C\beta_q$;
its bounds away from the center grow only polynomially with distance.
The weighted commutator calculus thus bounds $[D_t(f),Y(s)]$ by
$C\alpha_q\beta_q$.  Integrating
\[
 K_{q,d}(f)=\int_0^1 U_{q,d}(s)^*[D_t(f),Y(s)]U_{q,d}(s)\,ds
\]
proves the first assertion.  For the difference, use the same padded
$D(f)$ on both systems: its additional end-orbital number operators
commute with the smaller system's generator.  Thus no uncompensated
end value of $f$ occurs.  The source and generator
differences supply the factor $v^{-1}$ proved above; the difference of
the transporting evolutions supplies the same factor by Duhamel's
formula.  The weights $A_q$ and $v^{-1}$ have uniformly polynomial
ratios, so Lemma~\ref{locality:difference} applies also to their product.
All decompositions start with an annular source and expand its evolved
operator in shells about that source.  This gives the last assertion.
\end{proof}

The preceding norm estimate costs $C\int A_q(m)\,dm$.  Summation over
successive fluxes requires the stronger one-sided estimate below.
The weight $f$ is convex whereas the pinning logarithm $g$ is concave.
Their affine-subtracted quadratic parts therefore have opposite signs.
After filtering, these leading parts form a negative square; the cubic
Taylor remainder costs the additional factor $1/w$.

\begin{proposition}[One-sided change in a zero space]
\label{pinning:one-sided}
Writing $P_0$ for the zero projection of $H_t$, one has
\begin{equation}
 P_0K_{q,d}(f)P_0
 \le C\left(\int_0^t\frac{A_q(m)}{w(m)}\,dm\right)P_0.
 \label{eq:pinning-one-sided}
\end{equation}
\end{proposition}

\begin{proof}
Fix $s$ and abbreviate $H=H_{q,d}(s)$ and $P=P_{q,d}(s)$.
The derivative in the transported zero space is the compression of
\begin{equation}
 D_t(g)(1-P)D_t(f)+D_t(f)(1-P)D_t(g).
 \label{eq:pinning-compressed-derivative}
\end{equation}
We estimate this operator before integrating in $s$.

Choose a smooth $\chi_+$ equal to zero for
$\omega\le\gamma_{\mathrm{pin}}/2$ and to one for
$\omega\ge\gamma_{\mathrm{pin}}$.  Apply the filter
$r_+(\omega)=\chi_+(\omega)/\omega$ separately to the pieces in
Equation~\eqref{eq:pinning-source-decomposition}, and set
\[
 G_x=r_+(\operatorname{ad}_H)X_{g,x},\qquad
 F_x=r_+(\operatorname{ad}_H)X_{f,x}.
\]
As in the preceding filter argument, its time kernel has all absolute
polynomial moments.  Every piece has strictly positive spectral
frequencies.  Thus
\begin{equation}
 G_x^*P=F_x^*P=0,
 \qquad
 \left(\int G_x\,d\mu(x)\right)P=(1-P)D_t(g)P,
 \qquad
 \left(\int F_x\,d\mu(x)\right)P=(1-P)D_t(f)P.
 \label{eq:pinning-raising}
\end{equation}

At a center of latitude index $m$, put
$\rho_x=f''(m)/(-g''(m))>0$.  The common wavepacket decomposition gives
\begin{equation}
 F_x=-\rho_xG_x+E_x,
 \label{eq:pinning-local-proportionality}
\end{equation}
where $G_x$ has size $C\alpha_q(m)$ and $E_x$ has size
$C\beta_q(m)/w(m)$, at every fixed rapid order.  We verify the improved
error explicitly.  The function $a=f+\rho_x g$ has $a''(m)=0$.
Subtract its affine Taylor polynomial, as in
Equation~\eqref{eq:pinning-source-decomposition}.  In a row window of
width $w$, the third-order remainder is bounded, with all scaled
derivatives and polynomially weighted tails, by
\[
 Cw^3\left(\frac{f''(m)}{L+m}+
             \frac{\rho_x(-g''(m))}{v}\right)
 \le C\frac{\beta_q(m)}{w}.
\]
The last inequality uses $w^2\le u,v$ and $u\le L+m$.
For windows outside this central region, derivative ratios have
polynomial growth and the row tails absorb those powers.  The weighted
row estimate applied to $C_x(a)$ proves the claimed source error; the
same time filter preserves it.  This proves
Equation~\eqref{eq:pinning-local-proportionality} without making any
assertion about the sign of individual unfiltered row pieces.

Write $R_g=\int G_x\,d\mu(x)$, $R_f=\int F_x\,d\mu(x)$ and
$T=\int\sqrt{\rho_x}G_x\,d\mu(x)$.  Each integral is finite at fixed
flux.  From Equation~\eqref{eq:pinning-raising}, a reversed product with
a raising operator on the left vanishes in the compression, giving
$PG_x^*G_yP=P[G_x^*,G_y]P$.  Expanding
Equation~\eqref{eq:pinning-local-proportionality} therefore gives the
exact identity
\begin{align}
 P(R_g^*R_f+R_f^*R_g)P
 ={}&-2PT^*TP \notag\\
 &-\iint(\sqrt{\rho_x}-\sqrt{\rho_y})^2
          P[G_x^*,G_y]P\,d\mu(x)d\mu(y) \notag\\
 &+\iint P\bigl([G_x^*,E_y]+[E_x^*,G_y]\bigr)P
            \,d\mu(x)d\mu(y).
 \label{eq:pinning-negative-square}
\end{align}

We give the norm bounds for the last two integrals, since their spatial
summability is essential.  If $D=\operatorname{dist}(x,y)$, rapid
localization and the even parity of the pieces imply, for every fixed
$M$,
\[
 \|[G_x^*,G_y]\|\le C_M\alpha_x\alpha_y(1+D)^{-M},
 \qquad
 \|[G_x^*,E_y]\|\le
        C_M\alpha_x\frac{\beta_y}{w_y}(1+D)^{-M}.
\]
Here subscripts denote evaluation at the corresponding center.
The logarithmic derivative of $\rho$ is bounded by
$C((L+m)^{-1}+v^{-1})$.  Using radial resolution $w$ and the polynomial
ratio bounds of Section~\ref{sec:locality} consequently gives a fixed
exponent $k$ such that
\[
 |\sqrt{\rho_x}-\sqrt{\rho_y}|
 \le C\sqrt{\rho_x}\,\frac{D}{w_x}(1+D)^k.
\]
This bound follows by integrating the logarithmic derivative along the
meridian between the two latitudes when the distance is at most a small
multiple of $w_x$; outside that region its
right side dominates the polynomial ratio bound after increasing $k$.
Since $\rho_x\alpha_x^2\asymp A_x$, choose $M$ above all these
polynomial exponents plus the two-dimensional volume exponent.  Summing
first in $y$ bounds the second line of
Equation~\eqref{eq:pinning-negative-square} by
$C\int A_x/w_x^2\,d\mu(x)$, and its third line by
$C\int A_x/w_x\,d\mu(x)$.  Since $w\ge1$ and $-2PT^*TP\le0$,
Equation~\eqref{eq:pinning-compressed-derivative} is at most
$C\int_0^t A_q(m)/w(m)\,dm$ in quadratic form.  Conjugate back by
$U_{q,d}(s)$ and integrate over $0\le s\le1$ to obtain
Equation~\eqref{eq:pinning-one-sided}.
\end{proof}

\begin{corollary}[Cost of changing a calibration state]
\label{pinning:calibration-cost}
Let $q^0=3\lfloor q/3\rfloor$ and $r=q-q^0\in\{0,1,2\}$.
Consider the interaction $K_{q,3}(f)$ on $\mathcal F_{q-3}$ and its
common-frame comparison with $K_{q^0,3}(f)$.  Embed a unit vector
$\xi\in Z_{n,q^0-3}$ by the common orbital labels, or instead grow it
to flux $q-3$ using the $r$ isometries $\mathcal U_{Q,1}$.  The following
costs are each bounded by
\begin{equation}
 C\int_0^{q-3}\frac{A_q(m)}{v(m)}\,dm:
 \label{eq:pinning-calibration-cost}
\end{equation}
the change of interaction in the embedded state, and the change of its
expectation between the embedded and grown states.  The same bounds
hold after selecting any common collection of source-center pieces,
and after adding their absolute bounds.  In particular the statement
applies to the filled vector at flux $q^0-3$ and to separate northern
and southern portions of the interaction.
\end{corollary}

\begin{proof}
The first assertion is the integrated interaction difference in
Proposition~\ref{pinning:weighted-transport}; vacant extra modes are
compressed after comparison in the common frame.  For one growth step,
differentiate the expectation of the interaction along $U_{Q,1}(s)$.
The generator has size $\alpha_Q\le C/v$, while the observable
interaction has size $A_q$.  The weighted commutator calculus bounds
the derivative in operator norm by $C\int A_q/v$.  Finite-parameter
conjugation preserves this bound.  There are at most two growth steps,
and all size functions for their fluxes are comparable.  Applying this
argument to each selected source piece before summing proves the last
assertion, including comparisons with the directly embedded lower-flux
reference on common northern slots.
\end{proof}
